\section{Sonuç ve Gelecek Çalışmalar}
\label{CH:Sonuc}

Bu projede telefon kamerasıyla Türk İşaret Dili kelimelerini tanıyan, internetsiz
çalışan bir Android uygulaması geliştirilmiştir. AUTSL veri setinin işaretçiden
bağımsız bölmelerinde iki model karşılaştırılmış; 1B-CNN + Transformer modeli üç
tohumda yüzde $89{,}42 \pm 0{,}52$ test doğruluğuyla BiGRU'yu (yüzde
$87{,}98 \pm 0{,}23$) geride bırakmış ve daha az parametreyle yüzde 85 hedefini
aşmıştır. Model int8 LiteRT biçiminde 1,61 MB'tır ve dönüşüm doğruluğu
düşürmemiştir. Eğitim ve uygulama tarafında aynı MediaPipe modelleri ve birebir aynı
ön işleme kullanıldığı için telefon ve Python tahminleri neredeyse tamamen
örtüşmektedir.

Uygulama tez önerisindeki üç ekranı içermektedir. Bunlara ek olarak her kelimenin
örnek hareketini gösteren iskelet animasyonu ve gerçek ortam doğruluğunu ölçmek için
gönüllü testi ekranı eklenmiştir. Model boyutu ve AUTSL doğruluğu hedefleri
sağlanmış, internet bağlantısı gerektirmeyen bir yapı kurulmuştur. Uçtan uca hız ise
iki MediaPipe modelinin kare başına çalışması nedeniyle 15 FPS hedefinin altında
kalmıştır.

\subsection*{Gelecek çalışmalar}

\begin{itemize}
  \item \textbf{Canlı veriyle ince ayar:} Gönüllü testinde toplanan nokta kayıtları
        eğitim kümesine eklenerek model telefon kamerası koşullarına uyarlanabilir.
  \item \textbf{Sözlüğün genişletilmesi:} BosphorusSign22k (744 işaret) ve
        TurkSign446 (446 işaret) veri setlerine erişim izni alınarak sözlük
        büyütülebilir. TurkSign446 noktaları MediaPipe ile çıkarıldığı için mevcut
        hatta görece kolay eklenebilir.
  \item \textbf{Hız:} Analiz görüntüsünün küçültülmesi, el izlemenin önceki kareden
        başlatılması ve gövde modelinin daha seyrek çalıştırılması ile 15 FPS hedefine
        yaklaşılabilir.
  \item \textbf{Ayırt edici bilgi:} En çok karışan işaretler parmak teması ve derinlik
        bilgisiyle ayrışmaktadır. Elin derinlik koordinatı ve yüz noktalarının girdiye
        eklenmesi bu hataları azaltabilir.
  \item \textbf{Cümle düzeyi:} Ardışık kelimelerin dilbilgisel Türkçe cümleye
        çevrilmesi, sürekli işaret dili veri setleriyle ele alınabilir.
\end{itemize}
